\chapter{Bradycardia}

\section*{Definition and Overview}
Bradycardia is the medical term for a resting heart rate slower than normal, usually defined as fewer than 60 beats per minute in adults. For many healthy people, particularly young adults and well-trained athletes, a slow heart rate is a normal finding that simply reflects an efficient, fit cardiovascular system, and it requires no treatment. Bradycardia becomes a clinical problem when the heart beats too slowly to pump enough oxygen-rich blood to meet the body's needs, producing symptoms such as dizziness, fatigue, shortness of breath, and fainting. The severity ranges from an incidental finding on a routine check to a life-threatening rhythm disturbance, so the key question in every case is whether the slow rate is causing symptoms and whether an underlying cause needs treatment.

\section*{Epidemiology}
A slow resting heart rate is common in the general population, and its prevalence rises with age. It is seen more often in men than in women and is a normal finding in physically fit people, especially endurance athletes, in whom the resting rate can fall below 40 beats per minute without harm. The rhythm disorders that cause clinically significant bradycardia -- sick sinus syndrome and atrioventricular (AV) heart block -- increase markedly after the age of 60, and their incidence roughly doubles with each decade thereafter. Bradycardia is also common in people taking rate-lowering cardiac medications, in those with hypothyroidism, and in hospital settings after acute illness or heart attack.

\section*{Aetiology and Pathophysiology}
The heart rate is set by electrical impulses that begin in the heart's natural pacemaker, the sinoatrial node, and travel through the atria to the AV node and down the specialised conducting system to the ventricles. Bradycardia results when this electrical system fires too slowly, is suppressed, or is blocked along its pathway. Common causes include:
\begin{itemize}
    \item \textbf{Age-related degeneration:} progressive fibrosis of the sinoatrial node causes sick sinus syndrome, and fibrosis of the conducting system causes AV block of varying degrees.
    \item \textbf{Medications:} beta-blockers, non-dihydropyridine calcium-channel blockers, digoxin, and many anti-arrhythmic drugs slow the heart rate.
    \item \textbf{Thyroid disease:} an underactive thyroid gland (hypothyroidism) lowers metabolic demand and slows the heart.
    \item \textbf{Electrolyte imbalance:} particularly hyperkalaemia (high potassium), and less commonly hypocalcaemia or hypomagnesaemia.
    \item \textbf{Heart disease:} damage from myocardial infarction, myocarditis, cardiomyopathy, or coronary artery disease can disrupt the conducting system.
    \item \textbf{Other causes:} hypothermia, obstructive sleep apnoea, increased intracranial pressure, infections such as Lyme disease, and certain recreational drugs.
    \item \textbf{Physiological influences:} athletic conditioning, deep sleep, and vagal (vasovagal) reflexes all produce benign slowing of the heart rate.
\end{itemize}
Whatever the cause, the effect is the same: when the rate falls below the level needed to maintain adequate cardiac output, blood flow to the brain and other organs becomes insufficient, producing symptoms.

\section*{Clinical Features}
Many people with bradycardia have no symptoms at all, and the finding is incidental. When the heart rate is too slow to meet the body's needs, symptoms may include:
\begin{itemize}
    \item Dizziness, light-headedness, or a feeling of faintness
    \item Fainting or near-fainting (syncope), sometimes without warning
    \item Fatigue, weakness, and poor exercise tolerance
    \item Shortness of breath, especially with exertion
    \item Chest discomfort or a sensation that the heart is pausing or skipping beats
    \item Confusion, memory lapses, or difficulty concentrating, particularly in older people
    \item In severe cases, profound collapse, seizure, or cardiac arrest
\end{itemize}

\section*{Differential Diagnosis}
\begin{itemize}
    \item \textbf{Physiological bradycardia:} in athletes, during sleep, or with increased vagal tone -- normal and harmless.
    \item \textbf{Medication-induced bradycardia:} from beta-blockers, calcium-channel blockers, digoxin, and other rate-lowering drugs.
    \item \textbf{Sick sinus syndrome:} disease of the sinoatrial node in older adults, with alternating slow rates and episodes of tachycardia.
    \item \textbf{Atrioventricular (AV) block:} first, second, or third degree interruption of conduction between the atria and ventricles.
    \item \textbf{Hypothyroidism:} underactive thyroid gland causing generalised slowing, including of the heart.
    \item \textbf{Obstructive sleep apnoea:} nocturnal breathing pauses that stress the heart and cause daytime bradycardia.
    \item \textbf{Electrolyte disturbances and hypothermia:} high potassium or low body temperature suppress the heart rate.
    \item \textbf{Increased intracranial pressure:} from head injury or bleeding, which slows the heart via the nervous system.
\end{itemize}

\section*{When to Seek Medical Care}
See a doctor if you have an unusually slow pulse together with recurring dizziness, fainting, fatigue, shortness of breath, or chest discomfort, or if a prescribed heart medicine seems to be making you unwell. New symptoms, or symptoms that are worsening, should not be ignored, especially in older adults, in whom bradycardia may present as confusion or falls rather than as an obvious awareness of a slow pulse.

\textbf{Call 000 (or local emergency services) for:}
\begin{itemize}
    \item Fainting or collapse, especially without warning
    \item Chest pain, severe breathlessness, or difficulty breathing
    \item Confusion, extreme drowsiness, or difficulty waking
    \item A seizure or fitting episode
    \item A very slow pulse accompanied by any of the above
\end{itemize}

\section*{Diagnosis and Investigations}
\begin{itemize}
    \item \textbf{History and examination:} details of symptoms, medications, fitness level, and any underlying disease, together with measurement of the pulse, blood pressure, and heart sounds.
    \item \textbf{Electrocardiogram (ECG):} records the heart's electrical activity and identifies the specific rhythm disturbance, such as sinus bradycardia, sick sinus syndrome, or AV block.
    \item \textbf{Ambulatory monitoring:} a Holter monitor worn for 24--48 hours, or an event recorder worn for weeks, captures intermittent episodes and correlates symptoms with the rhythm.
    \item \textbf{Blood tests:} thyroid function, electrolytes, renal function, and cardiac markers to identify reversible causes.
    \item \textbf{Exercise testing:} reveals how the heart rate responds to activity and can unmask chronotropic incompetence.
    \item \textbf{Sleep study:} polysomnography if obstructive sleep apnoea is suspected.
    \item \textbf{Echocardiography:} assesses for structural heart disease when the cause is not otherwise apparent.
\end{itemize}

\section*{Management}
Treatment depends on whether the bradycardia is causing symptoms and on what is driving it:
\begin{itemize}
    \item \textbf{Observation:} no treatment is needed when the slow rate is physiological and symptom-free, as in athletes.
    \item \textbf{Medication review:} reducing or stopping rate-lowering drugs under medical supervision, and avoiding drugs that aggravate the condition.
    \item \textbf{Treating the underlying cause:} thyroid hormone replacement for hypothyroidism, treatment of sleep apnoea with CPAP, correction of electrolyte imbalance, and treatment of the relevant infection.
    \item \textbf{Pacemaker:} a small implanted device that senses the heart rate and paces it when it falls too low; it is highly effective and is the treatment of choice for symptomatic sick sinus syndrome and significant AV block.
    \item \textbf{Acute management:} in hospital, intravenous atropine or temporary pacing is used for unstable bradycardia while the cause is identified and corrected.
    \item \textbf{Follow-up:} regular review to monitor symptoms, adjust medication, and check the function of an implanted pacemaker.
\end{itemize}

\section*{Complications}
\begin{itemize}
    \item Recurrent fainting with the risk of falls, fractures, and head injury, particularly in older adults
    \item Reduced blood flow to the brain, causing confusion, and in severe cases collapse or impaired consciousness
    \item Reduced exercise tolerance and impaired quality of life from fatigue and breathlessness
    \item Rarely, high-grade or complete heart block with prolonged cardiac pauses that can lead to cardiac arrest
    \item Falls from unsteadiness, which themselves cause significant morbidity in the elderly
    \item Complications of pacemaker implantation, including infection, bleeding, and lead displacement, although these are uncommon
\end{itemize}

\section*{Prognosis and Outcomes}
For most people the outlook is excellent. Where bradycardia is physiological, as in athletes, no treatment is needed and it is not harmful. When it is caused by a reversible factor such as a medication or thyroid disease, correcting that factor usually resolves the problem completely. For persistent symptomatic bradycardia, a pacemaker is highly effective: most people feel significantly better after implantation, with dizziness and fainting resolving and exercise tolerance improving. Long-term studies show that people with appropriately treated symptomatic bradycardia have a good prognosis and return to normal activity.

\section*{Prevention and Self-Care}
\begin{itemize}
    \item Keep your heart healthy with regular physical activity, a balanced diet, and good control of blood pressure and cholesterol.
    \item Have prescribed heart medicines reviewed if they make you feel faint or unusually tired.
    \item Report fainting, dizzy spells, or unexplained fatigue to your doctor.
    \item Treat underlying conditions such as an underactive thyroid or sleep apnoea.
    \item Avoid excessive alcohol and recreational stimulant use, which can disturb heart rhythm.
    \item If you have a pacemaker, attend regular checks and carry your device information card.
    \item Do not stop or change rate-lowering medicines without medical advice.
\end{itemize}

\section*{Sources and Further Reading}
The following sources provide reliable, current information on bradycardia and its management.
\begin{itemize}
    \item NHS. \textit{Bradycardia.} https://www.nhs.uk/conditions/bradycardia/
    \item Mayo Clinic. \textit{Bradycardia: symptoms, causes, and treatment.} https://www.mayoclinic.org/
    \item Centers for Disease Control and Prevention. \textit{Heart disease and rhythm disorders.} https://www.cdc.gov/heartdisease/
    \item World Health Organization. \textit{Cardiovascular diseases.} https://www.who.int/health-topics/cardiovascular-diseases
    \item MSD Manual. \textit{Bradycardia.} https://www.msdmanuals.com/
    \item MedlinePlus. \textit{Bradycardia.} https://medlineplus.gov/
\end{itemize}
